% --- Archon physics formalization source begin ---
% archon:physics
% archon:covers PhyXMiniProblems/problem_phyx_mini_0340.lean
% archon:source-report reports/phyx_mini/problem_phyx_mini_0340.source.json

\chapter{Physics problem phyx\_mini\_0340}
\label{ch:PhyXMiniProblems_problem_phyx_mini_0340}

\paragraph{Problem source.}
\#\# Physical scenario

When a system is taken from state \$a\$ to state \$b\$ in \textbackslash{}textbf\{figure\}, 90.0\textasciitilde{}J of heat flows along path \$acb\$, and 60.0\textasciitilde{}J of work is done by the system.

\#\# Question

How much heat flows into the system along path \$adb\$ if the work done by the system is 15.0J?

\#\# Answer choices

- A: A: 15J

- B: B: 45J

- C: C: 25J

- D: D: 32J

\#\# Figure caption (auxiliary; use the image as primary evidence)

The image is a p-V (pressure-volume) diagram with labeled points and colored arrows indicating a thermodynamic cycle. 

- There are axes labeled \textbackslash{}( p \textbackslash{}) (pressure) on the vertical and \textbackslash{}( V \textbackslash{}) (volume) on the horizontal. The origin is labeled \textbackslash{}( O \textbackslash{}).
  
- Four points are marked: \textbackslash{}( a \textbackslash{}), \textbackslash{}( b \textbackslash{}), \textbackslash{}( c \textbackslash{}), and \textbackslash{}( d \textbackslash{}).

- The path consists of four segments forming a cycle:
  - From \textbackslash{}( a \textbackslash{}) to \textbackslash{}( b \textbackslash{}), there is a curved arrow (blue) representing an expansion process.
  - From \textbackslash{}( b \textbackslash{}) to \textbackslash{}( c \textbackslash{}), there is a horizontal arrow (purple) indicating an isochoric process (constant volume).
  - From \textbackslash{}( c \textbackslash{}) to \textbackslash{}( d \textbackslash{}), there is a vertical arrow (green) indicating a compression process.
  - From \textbackslash{}( d \textbackslash{}) to \textbackslash{}( a \textbackslash{}), there is a horizontal arrow (green) representing an isochoric process (constant volume).

- The arrows indicate the direction of the cycle, moving clockwise around the loop \textbackslash{}( a \textbackslash{}to b \textbackslash{}to c \textbackslash{}to d \textbackslash{}to a \textbackslash{}).

The diagram visually represents a cyclic process in a thermodynamic system.

\#\# Dataset metadata

Laws of Thermodynamics; Multi-Formula Reasoning, Physical Model Grounding Reasoning

\paragraph{Recorded answer/context.}
B: B: 45J

\paragraph{Figure/image path.}
/root/proposal\_for\_physic/science-mango/phyx\_mini\_run/phyx\_data/test\_image/340.png

\paragraph{Formalization target.}
create a compiling Lean file with sorry bodies at `PhyXMiniProblems/problem\_phyx\_mini\_0340.lean`. The Lean declarations must preserve the physical quantities, dimensions or dimensional roles, figure labels, governing-law hypotheses, and final relation expressed by this problem.
Use Mathlib/Physlib names found through LeanExplore where available. If a physics API is missing, introduce faithful local abstractions rather than scalar placeholder aliases.

\begin{theorem}[Physics formalization target]
\label{thm:physics:phyx_mini_0340:target}
\lean{PhyXMiniProblems.ProblemPhyXMini0340.heatIntoSystemAlongADB_eq_45_joules}
\uses{def:physics:phyx-mini-0340:phyxminiproblems-problemphyxmini0340-energyinjoules, def:physics:phyx-mini-0340:phyxminiproblems-problemphyxmini0340-pvfigure, def:physics:phyx-mini-0340:phyxminiproblems-problemphyxmini0340-matchessuppliedfigure, def:physics:phyx-mini-0340:phyxminiproblems-problemphyxmini0340-thermodynamicmodel, def:physics:phyx-mini-0340:phyxminiproblems-problemphyxmini0340-satisfiesfirstlaw}
For the two displayed paths from state \(a\) to state \(b\), if \(90\) joules
of heat enter and the system does \(60\) joules of work along \(acb\), while
the system does \(15\) joules of work along \(adb\), then \(45\) joules of
heat enter along \(adb\).
\end{theorem}
\begin{proof}
The figure hypotheses identify both \(acb\) and \(adb\) as paths with initial
state \(a\) and final state \(b\).  Applying
\(\Delta U=Q-W_{\mathrm{by}}\) to \(acb\) gives
\(\Delta U=90-60=30\) joules.  Internal energy depends only on the endpoint
state, so the same difference occurs along \(adb\).  A second application of
the first law gives \(30=Q_{adb}-15\), and hence \(Q_{adb}=45\) joules.
\end{proof}

% --- Archon declaration topology begin ---
\section{Declaration topology}
\begin{definition}[Volume Quantity]
  \label{def:physics:phyx-mini-0340:phyxminiproblems-problemphyxmini0340-volumequantity}
  \lean{PhyXMiniProblems.ProblemPhyXMini0340.VolumeQuantity}
  Physical volume, with dimension L³
\end{definition}
\begin{proof}
  This is the declared model definition of Volume Quantity.
\end{proof}
\begin{definition}[Pressure Quantity]
  \label{def:physics:phyx-mini-0340:phyxminiproblems-problemphyxmini0340-pressurequantity}
  \lean{PhyXMiniProblems.ProblemPhyXMini0340.PressureQuantity}
  Physical pressure, with dimension M L⁻¹ T⁻²
\end{definition}
\begin{proof}
  This is the declared model definition of Pressure Quantity.
\end{proof}
\begin{definition}[Energy Quantity]
  \label{def:physics:phyx-mini-0340:phyxminiproblems-problemphyxmini0340-energyquantity}
  \lean{PhyXMiniProblems.ProblemPhyXMini0340.EnergyQuantity}
  Heat, work, or internal energy, with dimension M L² T⁻²
\end{definition}
\begin{proof}
  This is the declared model definition of Energy Quantity.
\end{proof}
\begin{definition}[volume In Cubic Metres]
  \label{def:physics:phyx-mini-0340:phyxminiproblems-problemphyxmini0340-volumeincubicmetres}
  \lean{PhyXMiniProblems.ProblemPhyXMini0340.volumeInCubicMetres}
  \uses{def:physics:phyx-mini-0340:phyxminiproblems-problemphyxmini0340-volumequantity}
  SI cubic-metre readout of a physical volume
\end{definition}
\begin{proof}
  This is the declared model definition of volume In Cubic Metres.
\end{proof}
\begin{definition}[pressure In Pascals]
  \label{def:physics:phyx-mini-0340:phyxminiproblems-problemphyxmini0340-pressureinpascals}
  \lean{PhyXMiniProblems.ProblemPhyXMini0340.pressureInPascals}
  \uses{def:physics:phyx-mini-0340:phyxminiproblems-problemphyxmini0340-pressurequantity}
  SI pascal readout of a physical pressure
\end{definition}
\begin{proof}
  This is the declared model definition of pressure In Pascals.
\end{proof}
\begin{definition}[energy In Joules]
  \label{def:physics:phyx-mini-0340:phyxminiproblems-problemphyxmini0340-energyinjoules}
  \lean{PhyXMiniProblems.ProblemPhyXMini0340.energyInJoules}
  \uses{def:physics:phyx-mini-0340:phyxminiproblems-problemphyxmini0340-energyquantity}
  SI joule readout of a physical energy
\end{definition}
\begin{proof}
  This is the declared model definition of energy In Joules.
\end{proof}
\begin{definition}[Thermodynamic State]
  \label{def:physics:phyx-mini-0340:phyxminiproblems-problemphyxmini0340-thermodynamicstate}
  \lean{PhyXMiniProblems.ProblemPhyXMini0340.ThermodynamicState}
  \uses{def:physics:phyx-mini-0340:phyxminiproblems-problemphyxmini0340-volumequantity, def:physics:phyx-mini-0340:phyxminiproblems-problemphyxmini0340-pressurequantity}
  A point in the p-V state space of the system
\end{definition}
\begin{proof}
  This is the declared model definition of Thermodynamic State.
\end{proof}
\begin{definition}[Process Kind]
  \label{def:physics:phyx-mini-0340:phyxminiproblems-problemphyxmini0340-processkind}
  \lean{PhyXMiniProblems.ProblemPhyXMini0340.ProcessKind}
  Qualitative process type indicated by a segment in the p-V figure
\end{definition}
\begin{proof}
  This is the declared model definition of Process Kind.
\end{proof}
\begin{definition}[Process Leg]
  \label{def:physics:phyx-mini-0340:phyxminiproblems-problemphyxmini0340-processleg}
  \lean{PhyXMiniProblems.ProblemPhyXMini0340.ProcessLeg}
  \uses{def:physics:phyx-mini-0340:phyxminiproblems-problemphyxmini0340-thermodynamicstate, def:physics:phyx-mini-0340:phyxminiproblems-problemphyxmini0340-processkind}
  One directed leg of a thermodynamic path
\end{definition}
\begin{proof}
  This is the declared model definition of Process Leg.
\end{proof}
\begin{definition}[Thermodynamic Path]
  \label{def:physics:phyx-mini-0340:phyxminiproblems-problemphyxmini0340-thermodynamicpath}
  \lean{PhyXMiniProblems.ProblemPhyXMini0340.ThermodynamicPath}
  \uses{def:physics:phyx-mini-0340:phyxminiproblems-problemphyxmini0340-thermodynamicstate, def:physics:phyx-mini-0340:phyxminiproblems-problemphyxmini0340-processleg}
  A directed thermodynamic path, including its displayed legs
\end{definition}
\begin{proof}
  This is the declared model definition of Thermodynamic Path.
\end{proof}
\begin{definition}[Axis Quantity]
  \label{def:physics:phyx-mini-0340:phyxminiproblems-problemphyxmini0340-axisquantity}
  \lean{PhyXMiniProblems.ProblemPhyXMini0340.AxisQuantity}
  The physical quantity assigned to an axis in the supplied diagram
\end{definition}
\begin{proof}
  This is the declared model definition of Axis Quantity.
\end{proof}
\begin{definition}[PVFigure]
  \label{def:physics:phyx-mini-0340:phyxminiproblems-problemphyxmini0340-pvfigure}
  \lean{PhyXMiniProblems.ProblemPhyXMini0340.PVFigure}
  \uses{def:physics:phyx-mini-0340:phyxminiproblems-problemphyxmini0340-thermodynamicstate, def:physics:phyx-mini-0340:phyxminiproblems-problemphyxmini0340-thermodynamicpath, def:physics:phyx-mini-0340:phyxminiproblems-problemphyxmini0340-axisquantity}
  Raw labels, states, and directed paths read from the supplied p-V image
\end{definition}
\begin{proof}
  This is the declared model definition of PVFigure.
\end{proof}
\begin{definition}[Is Two Leg Path Via]
  \label{def:physics:phyx-mini-0340:phyxminiproblems-problemphyxmini0340-istwolegpathvia}
  \lean{PhyXMiniProblems.ProblemPhyXMini0340.IsTwoLegPathVia}
  \uses{def:physics:phyx-mini-0340:phyxminiproblems-problemphyxmini0340-thermodynamicstate, def:physics:phyx-mini-0340:phyxminiproblems-problemphyxmini0340-processkind, def:physics:phyx-mini-0340:phyxminiproblems-problemphyxmini0340-thermodynamicpath}
  A displayed two-leg path from startState through middleState to endState, with the indicated process type on each leg
\end{definition}
\begin{proof}
  This is the declared model definition of Is Two Leg Path Via.
\end{proof}
\begin{definition}[Is One Leg Path]
  \label{def:physics:phyx-mini-0340:phyxminiproblems-problemphyxmini0340-isonelegpath}
  \lean{PhyXMiniProblems.ProblemPhyXMini0340.IsOneLegPath}
  \uses{def:physics:phyx-mini-0340:phyxminiproblems-problemphyxmini0340-thermodynamicstate, def:physics:phyx-mini-0340:phyxminiproblems-problemphyxmini0340-processkind, def:physics:phyx-mini-0340:phyxminiproblems-problemphyxmini0340-thermodynamicpath}
  A displayed direct one-leg path with the indicated process type
\end{definition}
\begin{proof}
  This is the declared model definition of Is One Leg Path.
\end{proof}
\begin{definition}[Has Displayed PVGeometry]
  \label{def:physics:phyx-mini-0340:phyxminiproblems-problemphyxmini0340-hasdisplayedpvgeometry}
  \lean{PhyXMiniProblems.ProblemPhyXMini0340.HasDisplayedPVGeometry}
  \uses{def:physics:phyx-mini-0340:phyxminiproblems-problemphyxmini0340-volumeincubicmetres, def:physics:phyx-mini-0340:phyxminiproblems-problemphyxmini0340-pressureinpascals, def:physics:phyx-mini-0340:phyxminiproblems-problemphyxmini0340-thermodynamicstate}
  The rectangular placement of a, b, c, and d in the primary image. The strict SI inequalities record which pressure and volume levels are the lower and higher ones; the equalities record the isochoric and isobaric alignments
\end{definition}
\begin{proof}
  This is the declared model definition of Has Displayed PVGeometry.
\end{proof}
\begin{definition}[Matches Supplied Figure]
  \label{def:physics:phyx-mini-0340:phyxminiproblems-problemphyxmini0340-matchessuppliedfigure}
  \lean{PhyXMiniProblems.ProblemPhyXMini0340.MatchesSuppliedFigure}
  \uses{def:physics:phyx-mini-0340:phyxminiproblems-problemphyxmini0340-pvfigure, def:physics:phyx-mini-0340:phyxminiproblems-problemphyxmini0340-istwolegpathvia, def:physics:phyx-mini-0340:phyxminiproblems-problemphyxmini0340-isonelegpath, def:physics:phyx-mini-0340:phyxminiproblems-problemphyxmini0340-hasdisplayedpvgeometry}
  Exact axis, label, geometry, and arrow-direction information extracted from the supplied p-V image
\end{definition}
\begin{proof}
  This is the declared model definition of Matches Supplied Figure.
\end{proof}
\begin{definition}[Thermodynamic Model]
  \label{def:physics:phyx-mini-0340:phyxminiproblems-problemphyxmini0340-thermodynamicmodel}
  \lean{PhyXMiniProblems.ProblemPhyXMini0340.ThermodynamicModel}
  \uses{def:physics:phyx-mini-0340:phyxminiproblems-problemphyxmini0340-energyquantity, def:physics:phyx-mini-0340:phyxminiproblems-problemphyxmini0340-thermodynamicstate, def:physics:phyx-mini-0340:phyxminiproblems-problemphyxmini0340-thermodynamicpath}
  State-dependent internal energy and path-dependent heat and work. heatIntoSystem is positive for heat entering the system, while workDoneBySystem is positive for work done by the system
\end{definition}
\begin{proof}
  This is the declared model definition of Thermodynamic Model.
\end{proof}
\begin{definition}[Satisfies First Law]
  \label{def:physics:phyx-mini-0340:phyxminiproblems-problemphyxmini0340-satisfiesfirstlaw}
  \lean{PhyXMiniProblems.ProblemPhyXMini0340.SatisfiesFirstLaw}
  \uses{def:physics:phyx-mini-0340:phyxminiproblems-problemphyxmini0340-energyinjoules, def:physics:phyx-mini-0340:phyxminiproblems-problemphyxmini0340-thermodynamicpath, def:physics:phyx-mini-0340:phyxminiproblems-problemphyxmini0340-thermodynamicmodel}
  The first law with the sign convention ΔU = Q - W\_by, expressed in coherent SI joule readouts. Internal energy is a function of state, whereas heat and work are functions of the complete path
\end{definition}
\begin{proof}
  This is the declared model definition of Satisfies First Law.
\end{proof}
\begin{definition}[Answer Choice]
  \label{def:physics:phyx-mini-0340:phyxminiproblems-problemphyxmini0340-answerchoice}
  \lean{PhyXMiniProblems.ProblemPhyXMini0340.AnswerChoice}
  Answer labels printed with the problem
\end{definition}
\begin{proof}
  This is the declared model definition of Answer Choice.
\end{proof}
\begin{definition}[joule Value]
  \label{def:physics:phyx-mini-0340:phyxminiproblems-problemphyxmini0340-answerchoice-joulevalue}
  \lean{PhyXMiniProblems.ProblemPhyXMini0340.AnswerChoice.jouleValue}
  \uses{def:physics:phyx-mini-0340:phyxminiproblems-problemphyxmini0340-answerchoice}
  Joule value printed beside each answer label
\end{definition}
\begin{proof}
  This is the declared model definition of joule Value.
\end{proof}
% --- Archon declaration topology end ---
% --- Archon physics formalization source end ---
